\chapter{Electromagnetism and Fields}

\section{Maxwell's Equations and Mesh Networking}
In a fully decoupled sovereign node (Scenario Omega), communication cannot rely on centralized, legacy ISP infrastructure. Peer-to-peer (P2P) telemetry and ledger consensus are maintained via localized radio frequency (RF) meshes, typically utilizing LoRaWAN or WiFi-Direct protocols. The physical propagation of these signals is governed by Maxwell's equations, specifically the wave equation derived from Faraday's and Ampère's laws in free space:
\begin{equation}
\nabla^2 \mathbf{E} = \mu_0 \epsilon_0 \frac{\partial^2 \mathbf{E}}{\partial t^2}, \quad \nabla^2 \mathbf{B} = \mu_0 \epsilon_0 \frac{\partial^2 \mathbf{B}}{\partial t^2}
\end{equation}

To simulate the reliability of the Gossipsub network across the bioregion, the engine calculates the signal attenuation between any two nodes. In an idealized, unobstructed environment, the received power $P_r$ is dictated by the Friis transmission equation:
\begin{equation}
P_r = P_t G_t G_r \left( \frac{\lambda}{4\pi d} \right)^2
\end{equation}
where $P_t$ is transmitted power, $G_t$ and $G_r$ are antenna gains, $\lambda$ is the wavelength, and $d$ is the spatial distance \cite{balanis2015antenna}. However, the engine modifies this with material-specific attenuation coefficients when the line-of-sight intersects with high-density voxels (e.g., concrete walls, dense forest canopy), forcing nodes to strategically position repeaters to prevent network partition.

\section{Induction Electromagnetics and Smelting}
Electromagnetic interactions are not limited to data transmission; they are the primary vector for high-temperature material refinement (Scenario Theta: The Foundry). Electric induction kilns bypass the thermal resistance of traditional heating elements by directly inducing current within the conductive scrap metal.

According to Faraday's Law of Induction, an alternating magnetic field $\mathbf{B}$ generated by the kiln's water-cooled copper coils induces an electromotive force (EMF), driving eddy currents within the target metal:
\begin{equation}
\nabla \times \mathbf{E} = - \frac{\partial \mathbf{B}}{\partial t}
\end{equation}

These eddy currents encounter the electrical resistivity $\rho$ of the metal, generating immense, localized Joule heating ($P = I^2 R$). Because high-frequency alternating current is subject to the skin effect, the induced current is concentrated at the surface of the melt. The skin depth $\delta$ is critical for simulating the efficiency of the foundry:
\begin{equation}
\delta = \sqrt{\frac{2\rho}{\omega \mu}}
\end{equation}
where $\omega$ is the angular frequency of the AC current, and $\mu$ is the magnetic permeability of the metal \cite{davies1990induction}. The physics engine dynamically checks the $\delta$ of the loaded voxel entity (e.g., aluminum vs. iron) to determine the exact kW/h exergy draw required to achieve a phase transition.

\begin{thebibliography}{99}
\bibitem{balanis2015antenna}
Balanis, C. A. (2015). \textit{Antenna theory: analysis and design}. John Wiley \& Sons.

\bibitem{davies1990induction}
Davies, J. (1990). \textit{Induction heating handbook}. McGraw-Hill.

\end{thebibliography}
